\documentclass[11pt]{article}
\usepackage[a4paper,margin=26mm]{geometry}
\usepackage{amsmath,amssymb,amsthm,booktabs,hyperref}
\hypersetup{pdftitle={Boundary-sensitive products and periodic carry patterns in a two-dimensional radix system},pdfauthor={David Gaden}}
\newtheorem{theorem}{Theorem}[section]
\newtheorem{proposition}[theorem]{Proposition}
\newtheorem{corollary}[theorem]{Corollary}
\newtheorem{definition}[theorem]{Definition}
\newtheorem{remark}[theorem]{Remark}
\newcommand{\Z}{\mathbb Z}
\newcommand{\Q}{\mathbb Q}
\newcommand{\A}{\mathcal A}
\newcommand{\D}{\mathcal D}
\title{Boundary-sensitive products and periodic carry patterns\\in a two-dimensional radix system}
\author{David Gaden}
\date{12 October 2026}
\begin{document}
\maketitle
\begin{abstract}
We study finite integer arrays on a two-dimensional lattice under the
carry move that removes ten units at one site and adds one at each of its
two forward neighbours. A fixed signed digit convention gives a finite
canonical representative. Oppositely directed truncated fractional
cascades retain exact endpoint identities. For a product of $a\ge1$ cascades
in one direction and $c\ge1$ in the other, at a common cutoff, we give a finite
necessary-and-sufficient test for its eventual canonical display period.
The least period divides $10^{a+c-1}$ and is computable without enumerating
a full cycle. Limiting displays alone do not support multiplication.
As a geometric application, an injective two-channel encoding admits exact
dot-product and determinant tests with fixed nonzero array-valued factors.
Literal integer endpoint encoding loses directional information, whereas
retained directional units preserve it.
\end{abstract}

\noindent\textbf{Draft status.} This manuscript assembles current proofs and
reproducible experiments for specialist review. Priority and novelty have
not been established by a complete literature review. No physical model,
performance advantage, or complete arithmetic of arbitrary infinite real
inputs is claimed. The large phase-collision example remains in a separate
companion certificate and is not needed for the main theorem here.

\section{Context and elementary arithmetic}

The author developed the underlying idea around 2000 and explored it with
a C++ numerical application at that time. A new implementation followed
in 2025. The present work formalizes patterns observed in these experiments,
particularly the interaction of oppositely directed fractional cascades.

Multi-dimensional numbers (MDNs), called 2DNs in two dimensions, arrange
integer digits on a lattice. The entire spreadsheet-like grid represents
one number; its cells are the digits of that number. Positional order
increases rightward and upward, with degree $i+j$ at site $(i,j)$.
The horizontal convention is reversed from an ordinary written numeral.
This positional order does not assign an ordinary scalar weight to either
axis. Arithmetic is specified by the carry rule.

A finite grid has finitely many nonzero cells. Two finite grids represent
the same MDN when a finite sequence of carries connects them, allowing
negative as well as positive carry amounts. A carry of $k\in\Z$ at
$(i,j)$ makes the three changes
\[
 d_{i,j}\mapsto d_{i,j}-10k,\qquad
 d_{i+1,j}\mapsto d_{i+1,j}+k,\qquad
 d_{i,j+1}\mapsto d_{i,j+1}+k.
\]
Addition is coefficientwise, and multiplication is finite convolution.
They are associative, commutative and distributive modulo carries.

Laurent polynomials record positions, with $X$ and $Y$ denoting shifts.
For general integer base $b\ge3$, the finite carry algebra is
\[
 \A_b=\Z[X^{\pm1},Y^{\pm1}]/(X+Y-b).
\]
Here $X+Y=b$ is exactly the local carry, not an assignment of ordinary
scalar magnitudes to either axis. Eliminating the formal symbol $Y$ gives
\[
 \A_b\cong\Z[X,X^{-1},(b-X)^{-1}]\subset\Q(X).
\]
This algebraic embedding proves that the finite algebra has no zero
divisors and that $1\ne0$.

Digit systems in quotient rings have an established literature
\cite{digits}. The polynomial $10-X-Y$ is a sandpile polynomial under
Strasser's definition, and his local toppling move is the positive carry
above \cite{sandpile}. Signed digits, full finite boundaries, and
multiplication of directional cascade products require the specific
treatment given here. The finite-difference methods used below also belong
to the established study of modular binomial sequences \cite{binomials}.
These connections locate the problem; they do not establish priority for
the results of this draft.

\subsection{Polymorphism, signs, symmetry, and changing base}
Different grids can represent one carry class, a phenomenon termed
\emph{polymorphism} in the MDN interface. For example,
$1=-9+X+Y$ in base ten. Consequently the sign of an individual digit is
representation-dependent. Additive negation of a complete number is
well-defined, but the carry rule does not specify a total order.

In fact no total additive order on $\A_{10}$ can be invariant under
axis exchange $\sigma(X)=Y$, $\sigma(Y)=X$. The nonzero element $X-Y$
would have to have the same sign as its negative $Y-X$. Total orders
obtained by additional choices are not excluded by this argument.

Integer input at the origin is invariant under $\sigma$, and so is its
canonical display defined below. The converse is false: $XY$ is symmetric
but is not an ordinary integer. The unity line $i+j=0$ is different from
the reflection axis $i=j$. Selecting one fractional direction can break
axis-exchange symmetry; multiplying opposite directions can restore it.
Thus being below the unity line does not itself imply asymmetry.

Base is part of the carry relation. Keeping the same two shift labels
while changing $b$ cannot preserve both carries and the embedded integers:
the same $X+Y$ would equal two different integers. This does not imply
that every encoding between bases loses information. For example, after
extending coefficients to $\Q$, the substitution
$X_b\mapsto(b/c)X_c$, $Y_b\mapsto(b/c)Y_c$ gives an isomorphism between
the base-$b$ and base-$c$ algebras. It need not preserve the subset of finite
integer-coefficient grids. A claim about base conversion therefore needs
an explicit conversion map and coefficient domain.

\section{A canonical finite display}
Fix $\D=\{-4,-3,\ldots,5\}$ and define
\[
 h(v)=\left\lfloor\frac{v+4}{10}\right\rfloor,
 \qquad r(v)=v-10h(v)\in\D.
\]
Sweep through diagonals of increasing $i+j$, leaving $r(v)$ and forwarding
$h(v)$ to both neighbours.

\begin{proposition}
Every finite grid has exactly one finite representative with digits in
$\D$, and the sweep produces it.
\end{proposition}
\begin{proof}
Above the highest original input diagonal, if incoming coefficients have
absolute value at most $M$, the next layer has bound
$2\lfloor(M+5)/10\rfloor$. This is strictly smaller than $M$ for $M\ge2$;
for $M\le1$ no carry occurs. Each layer has finite support, so the sweep
terminates. For uniqueness, a finite zero-equivalent grid is
$(X+Y-10)H$ with finite $H$. Its lowest nonzero diagonal has all
coefficients divisible by ten. The difference of two $\D$-grids has every
coefficient between $-9$ and $9$, so a nonzero such difference cannot be
zero-equivalent.
\end{proof}

This normalization is a comparison convention. It is not an assertion
that the GUI's Positive sign convention is canonical.

\section{Directional cascades and retained boundaries}
Use positive term count $n$; the GUI's zero-based depth is $N=n-1$.
Define finite grids
\[
 L_n=\sum_{i=0}^{n-1}(-1)^iX^{-i-1}Y^i,
 \qquad R_n=\sum_{i=0}^{n-1}(-1)^iX^iY^{-i-1}.
\]
They are the finite directional entries labelled $0.1X$ and $0.1Y$.
Set $q=-Y/X$ and $E_n=q^n$. Telescoping gives
\begin{equation}\label{eq:boundary}
 10L_n=1-E_n,\qquad10R_n=1-E_n^{-1},\qquad
 10^{a+c}L_n^aR_n^c=(1-E_n)^a(1-E_n^{-1})^c.
\end{equation}
All endpoints are retained. In particular
\begin{equation}\label{eq:fuzzy}
 100L_nR_n=2-E_n-E_n^{-1},\qquad 10L_nR_n=L_n+R_n.
\end{equation}
The second identity follows from the first and torsion-freeness.

One possible rigorous infinite object is an eventual equivalence class of
sequences of complete finite grids: two sequences agree when their full
carry classes agree at every sufficiently late stage. Addition and
multiplication are stagewise and well-defined. This records an infinite
recipe; it does not assert that visible limiting arrays support those
operations. In such a recipe $E_nE_n^{-1}=1$ at every stage, so identifying
both endpoints with zero is incompatible with multiplication.

For independent cutoffs the exact identity instead is
\[
 10^{a+c}\prod_iL_{n_i}\prod_jR_{m_j}
 =\prod_i(1-q^{n_i})\prod_j(1-q^{-m_j}).
\]
The period theorem below concerns common cutoffs only.

\section{An exact algorithm for all matched-product periods}
Let $a,c\ge1$, $d=a+c$, $r=d-1$, and $G_n=L_n^aR_n^c$.
A shift $s>0$ is an \emph{eventual display period} if the canonical displays
of $G_{n+s}$ and $G_n$ agree in each fixed finite window for all sufficiently
large $n$. The threshold may depend on the window.

Before carrying, the coefficient at $(k-a,-k-c)$ is $(-1)^k C_n(k)$, with
\[
 C_n(k)=[z^{k+a(n-1)}](1+z+\cdots+z^{n-1})^d.
\]
Expansion as $(1-z^n)^d/(1-z)^d$ gives, for fixed $k$ and sufficiently
large $n$, $C_n(k)=P(n,k)+J(k)$, where
\[
 P(n,k)=\sum_{u=0}^{a-1}(-1)^u\binom du
                 \binom{(a-u)n+k+c-1}{r}
\]
and $J(k)$ is independent of $n$. Specifically, $J(k)=0$ for $k<a$ and
$J(k)=(-1)^a\binom da\binom{k+c-1}{r}$ otherwise. Binomials with arbitrary
integer upper arguments are interpreted as falling-product polynomials.

For $0\le j<r$, define
\begin{equation}\label{eq:Q}
 Q_{j,s}(n,k)=\sum_{u=0}^{a-1}(-1)^u\binom du
 \left[\binom{(a-u)(n+s)+k+c-1}{r-j}
       -\binom{(a-u)n+k+c-1}{r-j}\right].
\end{equation}

\begin{theorem}\label{thm:period}
A positive integer $s$ is an eventual display period precisely when
\begin{equation}\label{eq:finite-test}
 10^{j+1}\mid Q_{j,s}(n,k)
 \quad(0\le j<r,\ n,k\ge0,\ n+k\le r-j-1).
\end{equation}
The minimal period $T_{a,c}$ divides $10^r$. Starting with $10^r$ and
successively removing factors of 2 and 5 whenever \eqref{eq:finite-test}
passes gives exactly $T_{a,c}$.
\end{theorem}
\begin{proof}
Index layer $j$ by $(k-a,-k-c+j)$. Incoming coefficients obey
$v_0(k)=(-1)^kC_n(k)$ and
$v_{j+1}(k)=h(v_j(k))+h(v_j(k-1))$.
Equal incoming residues modulo ten give equal digits. Conditional on
agreement of previous layers, the incoming difference for a cutoff shift
$s$ is
\[
 \frac{(-1)^k}{10^j}\nabla_k^j\{P(n+s,k)-P(n,k)\},
 \qquad \nabla f(k)=f(k)-f(k-1).
\]
Forward differences equal backward differences after translating $k$,
and the $j$th forward difference is \eqref{eq:Q}. The initial difference
has degree at most $r-1$ in $k$, so its $r$th difference vanishes.
Thus agreement of all digits is equivalent to divisibility by $10^{j+1}$
for every integer $n,k$ and $0\le j<r$.

For completeness this equivalence also holds with eventual agreement.
Every fixed digit depends on only finitely many raw ancestors, putting
them in the polynomial interior for sufficiently large $n$. Conversely,
an integer-valued polynomial of degree at most $t$ is periodic modulo $M$
in each argument: a shift by $M\operatorname{lcm}(1,\ldots,t)$ preserves
its binomial-basis expansion. Eventual divisibility therefore implies
divisibility at every integer argument. Constant polynomials are immediate.

The total degree of $Q_{j,s}$ in $(n,k)$ is at most $t=r-j-1$.
Its Newton expansion has the form
\[
 Q(n,k)=\sum_{i+h\le t}(\Delta_n^i\Delta_k^hQ)(0,0)
                   \binom ni\binom kh.
\]
All coefficients are integer combinations of values in the finite triangle
$n,k\ge0$, $n+k\le t$. Divisibility there is therefore necessary and
sufficient everywhere, proving \eqref{eq:finite-test}.

Set $U=10^r$. In the binomial translation formula for each summand of
\eqref{eq:Q}, the coefficients include $\binom{vU}{h}$, where
$v=a-u\ge1$ and $1\le h\le r-j$. For $p=2,5$,
\[
 v_p\binom{vU}{h}\ge r-v_p(h)\ge r-(h-1)\ge j+1.
\]
Every summand is divisible by $10^{j+1}$, so $U$ is a period. Integer
periods form a subgroup of $\Z$ (negative shifts interpreted eventually),
so the least positive period divides $U$. A failed prime division means
its current prime exponent is minimal; removing the other prime cannot
change this. The stated reduction algorithm is exact.
\end{proof}

The finite test has at most $r(r+1)(r+2)/6$ sample congruences per candidate.
Its $n=0$ nodes evaluate the interior polynomial; they do not describe
a physical zero-term cascade.
If either factor count is zero, all fixed raw coefficients stabilize,
and carry locality gives period 1.

\begin{corollary}\label{cor:bound}
Writing $d=a+c$, the integer
\[
 T_0=\frac{10^{d-1}}{\gcd(10^{d-1},\binom{d-2}{a-1})}
\]
divides $T_{a,c}$. In particular $T_{1,c}=10^c$ and $T_{a,1}=10^a$.
\end{corollary}
\begin{proof}
The highest comparison difference is
$Q_{r-1,s}=(-1)^{a-1}\binom{d-2}{a-1}s$ by the alternating binomial sum.
Its required divisibility gives the lower bound. For $a=1$ or $c=1$
it coincides with the upper bound.
\end{proof}

The short lower bound need not be the period. For $(a,c)=(2,8)$ it gives
$125{,}000{,}000$, but at $j=6,n=1,k=0$ the remainder in
\eqref{eq:finite-test} is $5{,}000{,}000$ modulo $10^7$. The exact
algorithm gives $250{,}000{,}000$.

\begin{center}
\begin{tabular}{rrrrrrrr}
\toprule
$(a,c)$ &(1,1)&(2,1)&(2,2)&(3,1)&(3,2)&(4,2)&(3,3)\\
$T_{a,c}$ &10&100&500&1,000&10,000&25,000&50,000\\
\bottomrule
\end{tabular}
\end{center}

\section{Phases, hidden information, and a 10-adic coordinate}
The theorem produces a limiting array along any sequence of growing term
counts in a fixed period residue. These deterministic arrays are the
``fuzzy phases'': unrestricted growth need not select one of them.
Cycle length does not automatically count distinct arrays.

At a fixed residue label $n$, equality of phases at $n$ and $n+s$ is
equivalent to
\[
 10^{j+1}\mid Q_{j,s}(n,k),\qquad
 0\le j<r,\quad0\le k\le r-j-1.
\]
This is the same proof with a one-variable Newton test in $k$.

\begin{proposition}
Equality of limiting canonical arrays is not sufficient to define
multiplication consistently with these finite products.
\end{proposition}
\begin{proof}
Take term counts $n_t=100t+1$ and $m_t=n_t+10$. The $L_nR_n$ arrays
have the same limits, by period 10. The $L_n$ arrays also have the same
limit, by one-direction stabilization. Their products $L_n^2R_n$ have
different limits: the highest comparison difference in Corollary
\ref{cor:bound} requires a phase-label difference divisible by 100,
whereas $m_t-n_t=10$. Thus equal limiting inputs give unequal limiting
products.
\end{proof}

Compatible residues modulo $10,10^2,\ldots$ define a 10-adic integer
$\alpha\in\Z_{10}=\varprojlim\Z/10^k\Z$. Theorem \ref{thm:period}
shows that $\alpha$ determines the joint phase data for every matched
product: a product of degree $d$ needs at most its residue modulo
$10^{d-1}$. Conversely the probes $L_nR_n^c$ distinguish every residue
modulo $10^c$, since their highest comparison difference is simply $s$.
Thus the joint phase family distinguishes every such $\alpha$.
This is a coordinate on the joint observations, not an identification of
the complete boundary-retaining algebra with $\Z_{10}$.

\section{Vector encoding and a boundary-valued circle identity}
Fix $n\ge1$ and abbreviate $L=L_n$, $R=R_n$, $F=LR$, $E=q^n$.
The coefficient identities below may be read in $\A_{10}\otimes\Q$;
integer coefficient inputs use only the original finite-grid algebra.
The three elements $1,L,R$ are linearly independent over $\Q$: multiplying
$A+BL+CR=0$ by ten gives
$(10A+B+C)-BE-CE^{-1}=0$, forcing all coefficients to vanish.
Here $q$ and hence $E$ are transcendental formal elements.

Consequently $C(a,b)=aL+bR$ is an injective additive encoding of integer
vectors. It extends to rational vectors in the rational coefficient
extension. It is not a ring embedding of coordinatewise vector arithmetic.
The integer input parser gives the same element $1$ for literal $1X$ and
$1Y$, losing the distinction between $(1,0)$ and $(0,1)$.
Unlike the polar origin, these are different nonzero vectors.

Direction-preserving units instead are
\[
 U=10L=1-E,\qquad V=10R=1-E^{-1}.
\]
They are distinct and satisfy $UV=U+V$. No unital homomorphism to a
nonzero ring can send both to 1, since it would imply $1=2$.

For an exact unit-circle point $(s,t)=(3/5,4/5)$ the native decimal
inputs are $6L$ and $8R$. The two proposed circle expressions are
\[
 (6L)(6R)+(8L)(8R)=100F=2-E-E^{-1}\ne1,
\]
\[
 100\bigl((6L)^2+(8R)^2\bigr)
 =36(1-E)^2+64(1-E^{-1})^2\ne100.
\]
For any rational $s,t$ with $s^2+t^2=1$, the consistent linear
directional encoding gives the useful constant identity
\[
 (sU)(sV)+(tU)(tV)=UV.
\]
The constant retains both endpoints. This extends to formal scalars with
$s^2+t^2=1$; using arbitrary real trigonometric scalars is an additional
coefficient extension, not a proved property of arbitrary decimal parsing.

\begin{proposition}[Dot product and oriented area]
Let $C(a,b)=aL+bR$, let $S C(a,b)=aL-bR$ on this encoded vector space,
and let $\sigma$ exchange axes. Then
\begin{align*}
 C(a,b)\sigma C(c,d)+S C(a,b)\sigma S C(c,d)
   &=2(ac+bd)F,\\
 C(a,b)\sigma C(c,d)-\sigma C(a,b)C(c,d)
   &=(ad-bc)(L^2-R^2).
\end{align*}
The first expression vanishes exactly for orthogonal vectors; the second
vanishes exactly for linearly dependent vectors.
\end{proposition}
\begin{proof}
Expand by ordinary distributivity. In the first sum the $L^2$ and $R^2$
terms cancel; in the second difference the $LR$ terms cancel. The factors
$F$ and $L^2-R^2$ are nonzero in the domain, so the zero tests are exact.
\end{proof}

Thus selected combinations of products implement geometric tests while
retaining every boundary. The full product always has all four
distributive terms. The operation $S$ is defined on the vector subspace;
it is not asserted to be a ring automorphism. Similarly a quarter-turn
$J C(a,b)=C(-b,a)$ satisfies $J^2=-I$ on that subspace. This does not
identify an MDN axis with an imaginary unit or turn native multiplication
into complex multiplication.

For $(a,b)=(6,8)$ and $(c,d)=(-8,6)$, the determinant expression is
\[
100(L^2-R^2)=-2E+E^2+2E^{-1}-E^{-2}.
\]
All of this nonzero area information is at the retained endpoints.
A sufficiently small central window therefore misses it completely.

\section{Computational evidence and remaining work}
The repository's \texttt{research/general\_products} contains the exact
finite-test implementation and 120 minimal-period certificates through
total degree sixteen. A native probe linked the corrected MDN library and
checked 989 finite cases with 14,832 exact checks and no failures.
These compare full canonical grids and explicit carry certificates,
including boundary identities. Separate exact local formula calculations
check very large cutoff shifts; full native grids at those huge depths
were not constructed.

The new \texttt{research/vector\_circle} probe passed 45 cases and 4,045
exact checks, with zero failures. It tests rational unit-circle
inputs, literal axis endpoints, full distributivity, and paired geometric
identities at finite depths. Its generated \texttt{backend\_results.json}
is the authoritative record of the current test count and outcomes.
Both studies use the corrected C++ library. The
\href{https://github.com/Marupio/mdn-research}{repository} includes the
backend corrections, regression tests, research scripts and compact result
records. Reproduction commands are collected in
\texttt{research/manuscript/README.md}; full generated carry certificates
can be regenerated. A permanent archival release and identifier remain
necessary before submission.

The next mathematical problems are to characterize a minimal
multiplication-compatible boundary state, extend period analysis to
independent cutoff paths, and determine the relation of these results to
existing radix and algebraic-dynamics theorems. Exact geometric identities
alone establish neither a faster numerical method nor a physical model.
Independent specialist review and an explicit novelty comparison are
required before a submission claim is made.

\paragraph{Development and assistance.}
The original construction, C++ application and motivating observations
predate modern large language models. Recent AI assistance helped translate
the author's computational descriptions and conjectures into mathematical
definitions, proof drafts and reproducible checks. This assistance does not
constitute independent specialist review.

\begin{thebibliography}{9}
\bibitem{digits} K. Scheicher, P. Surer, J. M. Thuswaldner and
C. E. van de Woestijne, \emph{Digit systems over commutative rings},
\href{https://arxiv.org/abs/1004.3729}{arXiv:1004.3729}.
\bibitem{sandpile} G. Strasser, \emph{Abelian sandpiles and algebraic
models}, \href{https://arxiv.org/abs/1707.02281}{arXiv:1707.02281}.
\bibitem{binomials} L. Fiorot, R. Gilblas and A. Tonolo,
\emph{Periodic sequences, binomials modulo a prime power, and a math/music
application}, Advances in Applied Mathematics 162 (2025), 102786,
\href{https://doi.org/10.1016/j.aam.2024.102786}{doi:10.1016/j.aam.2024.102786}.
\end{thebibliography}
\end{document}
